\documentclass[11pt]{article}
\usepackage[margin=1in]{geometry}
\usepackage{amsmath,amssymb,amsthm}
\usepackage{graphicx}
\usepackage{microtype}
\usepackage[hidelinks]{hyperref}

\newtheorem{theorem}{Theorem}
\newtheorem{proposition}{Proposition}
\newtheorem{remark}{Remark}
\newcommand{\Hh}{\mathcal H}
\newcommand{\R}{\mathbb R}

\title{A continuous Tauberian theorem for Karcher means\\
and a sharp obstruction to the unit-lag formulation}
\author{Candidate full answer for arXiv:1906.06988}
\date{23 July 2026}

\begin{document}
\maketitle

\begin{center}
\fbox{\parbox{0.91\textwidth}{
\textbf{Status: candidate full answer, likely valid.}
The curve analogue of the source's discrete Tauberian theorem holds in every
Hadamard space under the natural metric-derivative condition.  An explicit
smooth real-valued curve shows that the superficially literal unit-time-lag
condition is insufficient.}}
\end{center}

\section{Source question}

The source is Hadi Khatibzadeh and Hadi Pouladi,
\emph{Tauberian Conditions for Almost Convergence in a Geodesic Metric
Space}, arXiv:1906.06988 \cite{KP}.  For a sequence $(x_n)$ in a Hadamard
space, their Theorem 3.5 says that convergence of the Karcher means
$\sigma_n$ to $y$, together with
\[
 n\,d(x_n,x_{n-1})\longrightarrow0,
\]
implies $x_n\to y$.  Remark 3.1, PDF page 14, defines the Karcher mean
$\sigma_T$ of a curve $c:[0,\infty)\to\Hh$ as the unique minimizer of
\[
 \mathcal G_T(z):=\frac1T\int_0^T d^2(c(t),z)\,dt
\]
and states that ``Theorem 3.5 for curves remains an open question.''

\clearpage
\begin{center}
\centering
\includegraphics[width=\textwidth]{figures/open_problem_crop.png}
\smallskip

{\small Figure 1. Remark 3.1 and the open curve question, source PDF page 14.}
\end{center}

The source does not spell out which continuous replacement for the discrete
increment condition is intended.  The result below gives the infinitesimal
analogue and also settles the tempting unit-lag interpretation negatively.

\section{Definitions and proof intuition}

A nonnegative $g\in L^1_{\mathrm{loc}}([0,\infty))$ is an
\emph{upper gradient along $c$} if
\[
 d(c(s),c(t))\leq\int_s^t g(u)\,du
 \qquad(0\leq s\leq t).
\]
Every locally absolutely continuous metric-space curve admits such a $g$;
one may take its metric derivative $g=|c'|$ almost everywhere.

\subsection*{Proof intuition}

The key observation is a first-moment estimate for a CAT(0) barycenter:
its distance from any comparison point is no greater than the average
distance of the data from that point.  Apply it with the comparison point
$c(T)$.  The upper-gradient inequality bounds the distance from $c(t)$ to
$c(T)$ by the variation accumulated between $t$ and $T$.  Reversing the
order of integration counts the speed at time $s$ with weight $s$, producing
exactly the continuous counterpart of the discrete Abel-summation term.

The unit-lag condition misses all motion inside each interval of length one.
A periodic sine curve therefore has zero unit-lag increments while retaining
nonvanishing oscillation.

\section{The continuous Tauberian theorem}

\begin{theorem}[Curve analogue of Theorem 3.5]\label{thm:main}
Let $(\Hh,d)$ be a Hadamard space and let
$c:[0,\infty)\to\Hh$ be locally absolutely continuous.  Let
$g\in L^1_{\mathrm{loc}}([0,\infty))$ be a nonnegative upper gradient along
$c$, and let $\sigma_T$ be the Karcher mean of $c|_{[0,T]}$.  Then for every
$y\in\Hh$ and $T>0$,
\begin{equation}\label{eq:quantitative}
 d(c(T),y)\leq d(\sigma_T,y)
 +\frac1T\int_0^T t\,g(t)\,dt.
\end{equation}
Consequently, if
\begin{equation}\label{eq:tauberian}
 \sigma_T\longrightarrow y
 \quad\hbox{and}\quad
 \frac1T\int_0^T t\,g(t)\,dt\longrightarrow0,
\end{equation}
then $c(T)\to y$.  In particular, the conclusion holds if
$t\,|c'|(t)\to0$ in the essential sense.
\end{theorem}

\begin{proof}
We first prove the barycenter estimate
\begin{equation}\label{eq:barycenter}
 d(\sigma_T,z)\leq\frac1T\int_0^T d(c(t),z)\,dt
 \qquad(z\in\Hh).
\end{equation}
Put $b=\sigma_T$, $D=d(b,z)$, and
\[
 F(w)=\frac1T\int_0^T d^2(c(t),w)\,dt.
\]
The CAT(0) strong-convexity inequality, applied on the geodesic from $b$ to
$z$, and the minimality of $b$ imply the variance inequality
\begin{equation}\label{eq:variance}
 D^2\leq F(z)-F(b).
\end{equation}
For each $t$, the triangle inequality gives
\[
 d^2(c(t),b)\geq
 \bigl(D-d(c(t),z)\bigr)^2,
\]
and hence
\[
 d^2(c(t),z)-d^2(c(t),b)
 \leq 2D\,d(c(t),z)-D^2.
\]
After integration and use of \eqref{eq:variance},
\[
 D^2\leq F(z)-F(b)
 \leq\frac{2D}{T}\int_0^T d(c(t),z)\,dt-D^2.
\]
If $D=0$, \eqref{eq:barycenter} is immediate; otherwise division by $2D$
proves it.

Now set $z=c(T)$ in \eqref{eq:barycenter}.  The upper-gradient property and
Tonelli's theorem give
\begin{align*}
 d(\sigma_T,c(T))
 &\leq\frac1T\int_0^T d(c(t),c(T))\,dt\\
 &\leq\frac1T\int_0^T\int_t^T g(s)\,ds\,dt\\
 &=\frac1T\int_0^T s\,g(s)\,ds.
\end{align*}
The triangle inequality proves \eqref{eq:quantitative}, and
\eqref{eq:tauberian} proves the convergence assertion.

Finally, if $t|c'|(t)\to0$ essentially, then its Cesaro averages tend to
zero: split the integral at a fixed sufficiently large time, divide the
fixed initial integral by $T$, and bound the remaining average by an
arbitrary $\varepsilon$.  Taking $g=|c'|$ completes the last assertion.
\end{proof}

\begin{remark}[A slightly weaker exact hypothesis]
The averaged condition in \eqref{eq:tauberian}, rather than pointwise decay
of $t|c'|(t)$, is all the proof requires.  This is the direct continuous
counterpart of the Cesaro expression that appears after Abel summation in the
source proof of Theorem 3.5.
\end{remark}

\section{The unit-lag interpretation is false}

\begin{proposition}\label{prop:counterexample}
The hypotheses
\[
 \sigma_T\longrightarrow y,
 \qquad
 T\,d(c(T),c(T-1))\longrightarrow0
\]
do not imply $c(T)\to y$, even for a bounded smooth curve in the Hadamard
space $\R$.
\end{proposition}

\begin{proof}
Take
\[
 c(t)=\sin(2\pi t),\qquad t\geq0.
\]
In $\R$, the Karcher mean for squared distance is the ordinary integral
mean, so
\[
 \sigma_T=\frac1T\int_0^T\sin(2\pi t)\,dt
 =\frac{1-\cos(2\pi T)}{2\pi T}\longrightarrow0.
\]
Periodicity gives $c(T)=c(T-1)$ for every $T\geq1$, so the unit-lag
Tauberian expression vanishes identically.  Nevertheless,
\[
 c(n+\tfrac14)=1,\qquad c(n+\tfrac34)=-1,
\]
and the curve does not converge.
\end{proof}

\section{Verification, scope, and novelty}

The argument is deterministic and has no computational dependency.  It uses
only CAT(0) strong convexity, minimality of the Karcher mean, the triangle
inequality, the upper-gradient property, and Tonelli's theorem.  The cases
$d(\sigma_T,z)=0$ and $d(\sigma_T,z)>0$ are separated before division.
No global boundedness assumption on $c$ is used.

Theorem~\ref{thm:main} answers the source's curve question for the natural
metric-derivative analogue in every Hadamard space.
Proposition~\ref{prop:counterexample} rules out the other immediate
interpretation.  The packet does not answer the source's separate questions
about uniformly convex non-CAT(0) spaces or about $\triangle$-almost
convergence.

The run's lightweight indexes were searched using the arXiv id, exact title,
and the terms \emph{almost convergence}, \emph{Karcher mean},
\emph{curve}, and \emph{Tauberian}.  A bounded web/arXiv search through
23 July 2026 used the exact title, the exact phrase ``Theorem 3.5 for
curves'', the phrase ``Karcher mean curve Tauberian condition'', and the
authors' names.  It found the source paper and later work using the same
curve-mean framework, but no resolution of this question.  This is not an
exhaustive novelty certification.

\section{Human-review recommendation}

The packet should be reviewed by an analyst familiar with CAT(0)
barycenters and Tauberian theory.  The mathematical proof is short; the main
interpretive question is whether the source authors intended the
metric-derivative or unit-lag translation.  Both possibilities are resolved
here.

\begin{thebibliography}{9}
\bibitem{KP}
H. Khatibzadeh and H. Pouladi,
\emph{Tauberian Conditions for Almost Convergence in a Geodesic Metric
Space}, arXiv:1906.06988 (2019).
\end{thebibliography}

\end{document}
